\documentclass[10pt, letterpaper]{article}

\usepackage[T1]{fontenc}
\usepackage[margin=0.5in]{geometry}

% Remove page numbers
\pagestyle{empty}

% Reduce spacing
\setlength{\parindent}{0pt}
\setlength{\parskip}{0pt}

% Custom section command
\newcommand{\ressection}[1]{%
    \vspace{10pt}%
    {\large\bfseries\MakeUppercase{#1}}%
    \vspace{2pt}%
    \hrule%
    \vspace{6pt}%
}

% Custom entry command
\newcommand{\resumeentry}[4]{%
    \textbf{#1}\hfill #2\\%
    \textit{#3}\hfill #4%
}

\begin{document}

% Header
\begin{center}
    {\LARGE\bfseries Patrick Manning Kennedy}\\[4pt]
    US Citizen \quad$\bullet$\quad 727-278-3108 \quad$\bullet$\quad
    pmkennedy4@gmail.com \quad$\bullet$\quad
    github.com/patrick4k
\end{center}

\ressection{Education}
\resumeentry{Embry-Riddle Aeronautical University (ERAU)}{Daytona Beach, FL}{Major: Aerospace Engineering (Astronautics), Minor: Computer Science, GPA: 3.943}{Dec 2024}

\ressection{Experience}

\vspace{4pt}
\resumeentry{ASUS Robotics \& AI Center}{}{Engineer}{Feb. 2025 -- Present}
\vspace{-6pt}
\begin{itemize}
    \setlength{\itemsep}{0pt}
    \setlength{\parskip}{0pt}
    \setlength{\topsep}{0pt}
    \item Architected a distributed, asynchronous messaging framework using Rust and C++ for dynamic autonomous systems.
    \item Researched, designed and implemented an efficient VIO system for real-time localization and control.
\end{itemize}


\resumeentry{NASA funded Swarm Unmanned Aerial Vehicle using Emergence (SUAVE), EPPL}{}{Software/Tech Lead}{Jan. 2024 -- Dec. 2024}
\vspace{-6pt}
\begin{itemize}
    \setlength{\itemsep}{0pt}
    \setlength{\parskip}{0pt}
    \setlength{\topsep}{0pt}
    \item Utilized \$86K budget to design and implement a GPS-denied UAV swarm for exploration and mapping.
    \item Research computer vision techniques including Visual-SLAM, HSV masking, and Collaborative SLAM.
    \item Computed local odometry with V-SLAM to enable GPS-denied flight using ROS2 and MAVLink in C++.
    \item Independently designed a robust drone controller interface in C++ with an emphasis on safety and portability.
    \item Implemented and tuned 4 PID controllers for a constrained ``follow me'' UAV position control.
    \item Developing a dual quaternion swarm controller that emulates dynamics and kinematics of a virtual structure to navigate a swarm of drones from one pose to another.
\end{itemize}

\vspace{4pt}
\resumeentry{Guava Compiler}{}{Independent Project}{Aug. 2023 -- Dec. 2024}
\vspace{-6pt}
\begin{itemize}
    \setlength{\itemsep}{0pt}
    \setlength{\parskip}{0pt}
    \setlength{\topsep}{0pt}
    \item Developing a statically typed language in C++ intended to generate safe and modular code.
    \item Designed a dynamic grammar focusing on readability while promoting modern design patterns and practices.
    \item Implementing macros in native Guava syntax to modify the intermediate representation (IR) on compile time.
    \item Exposing Guava compiler plugin points for custom runtime targets such as LLVM IR, transpiled GPGPU code, or interpretation, while maintaining interoperability between targets.
\end{itemize}

\vspace{4pt}
\resumeentry{Ansys Government Initiatives}{Exton, PA}{STK Communications \& Radar Developer}{May. 2023 -- Aug. 2024}
\vspace{-6pt}
\begin{itemize}
    \setlength{\itemsep}{0pt}
    \setlength{\parskip}{0pt}
    \setlength{\topsep}{0pt}
    \item Designed and developed complex radar workflows in C++ including 3D antenna gain pattern graphics, and real-time SAR radar image computation and overlay.
    \item Researched and interfaced with propagation models including GPU accelerated shooting and bouncing rays.
    \item Utilized COM and ATL in C++ and C\# for interoperability between different programs and API generation.
    \item Implemented gRPC to create a thread safe server client interface into STK, enabling a scenario dependent controller to be executed in parallel to a real-time simulation.
\end{itemize}

\vspace{4pt}
\resumeentry{Ansys Government Initiatives}{Exton, PA}{STK Quality Assurance Intern}{Sept. 2022 -- Dec. 2022}
\vspace{-6pt}
\begin{itemize}
    \setlength{\itemsep}{0pt}
    \setlength{\parskip}{0pt}
    \setlength{\topsep}{0pt}
    \item Conducted thorough vulnerability analysis of the 10-dimensional MODTRAN-derived model in Python.
    \item Developed regression test in Perl flagging interpolation inaccuracies within the MODTRAN-derived model.
    \item Evaluated EOIR Sensor API feature to identify release critical bugs or unexpected behaviors.
    \item Developed regression test in Perl ensuring functionality is retained for future releases across all edge cases.
\end{itemize}

% \vspace{4pt}
% \resumeentry{Phantom: Solid Motor Rocket, ERPL}{}{Project Lead}{May 2021 -- Aug. 2022}\\
% \textit{Propulsion Sub Team Lead}\hfill Sept. 2020 -- May 2021
% \vspace{-6pt}
% \begin{itemize}
%     \setlength{\itemsep}{0pt}
%     \setlength{\parskip}{0pt}
%     \setlength{\topsep}{0pt}
%     \item Improved upon the manufacturing process of complex solid rocket motor geometries.
%     \item Led team of six to build an internal solid rocket motor ballistics simulation in MATLAB and Java using an object-oriented design, resulting in thrust data within 9\% of market simulations.
% \end{itemize}

\ressection{Publications}
D. Golan, \textbf{P. Kennedy}, B. Gonzalez, et al., ``Swarm UAVs for area mapping in GPS-denied locations'', \textit{SPIE Defense and Commercial Sensing Conference Proceedings}, Paper No. 13055-28, June 2024.

\ressection{Relevant Skills}
\textbf{Languages:} Bash, C\#, C++, Java, JavaScript, Kotlin, MATLAB, Perl, Python, Rust\\[2pt]
\textbf{Programming:} ANTLR4, CMAKE, COM, Gazebo, Git, gRPC, Linux, .NET, OpenCV, Perforce, Regex, ROS2

\end{document}
